\documentclass[11pt,a4paper]{article}
\usepackage[T1]{fontenc}
\usepackage[utf8]{inputenc}
\usepackage{lmodern,microtype,amsmath,amssymb,booktabs}
\usepackage[margin=27mm]{geometry}
\usepackage[colorlinks=true,urlcolor=blue]{hyperref}
\hypersetup{pdftitle={Revision letter: Grothendieck rings of complex fusion categories, ranks 2--4},pdfauthor={Sébastien Palcoux}}
\setlength{\parindent}{0pt}
\setlength{\parskip}{6pt plus 1pt minus 1pt}
\raggedbottom
\begin{document}
{\small 19 November 2026}\par
{\Large\bfseries Revision letter}\par
{\large\itshape Grothendieck rings of complex fusion categories, ranks 2--4}\par
\medskip
Dear readers and moderators of Hexagon,

The three accompanying manuscripts are separate submissions, one for each rank.
They give a modern treatment of the rings of simple objects and their tensor products.
Each paper can be read on its own. The existence of a spherical structure is proved for
every categorification in the rank under study. The classification is up to
based-ring isomorphism, and the construction references provide examples realising the rings.

\begin{center}
\begin{tabular}{@{}lrrl@{}}\toprule
Paper & Previous pages & Revised pages & Based rings \\ \midrule
Rank two & 5 & 3 & 2\\
Rank three & 12 & 10 & 5\\
Rank four & 90 & 50 & 15\\ \bottomrule
\end{tabular}
\end{center}

\textbf{Rank two.}
For $X^2=1+nX$, let $m$ be the signed lift multiplicity. Integrality of the
global-dimension quotient of the sphericalization implies that
$(m^2+4)/(n^2+4)$ is a rational algebraic integer in $(0,1]$. This forces
$|m|=n$ and gives a spherical section. The induction matrix then leaves two
central twists whose sum is $-n\dim(X)$. Taking an algebraic conjugate gives
$n\operatorname{FPdim}(X)\le2$. This removes the earlier compositum calculation
and the separate Gauss-sum exclusion. The paper credits Ostrik for the
classification and identification of the central objects, Moore--Seiberg
for the surviving tensor structures, Thornton for near-group sphericality,
and Morrison--Walker for the induction-matrix method.

\textbf{Rank three.}
Both exclusions use the mass at twist one in a positive measure on the unit
circle. In the cubic case six even moments force the candidate and yield
spherical descent. In the spherical case the polynomial $3+4z+3z^2$
combines the first two moments in one inequality. It treats both the mixed
fusion rules and the near-group family, including the narrow endpoint
$m=3$, where the contradiction is $9\sqrt{17}\le37$. Consequently the
paper no longer takes Ostrik's pivotal structural reduction as an input.
It cites his quadratic descent theorem with its exact hypothesis.
The revised proof explains why the dimension character can be transported
to the positive embedding. The history distinguishes Bion-Nadal's
$E_6$ subfactor construction from Hagge--Hong's classification of the
corresponding tensor structures.

\newpage
\textbf{Rank four.}
The revised paper proves sphericality before applying the pivotal
classification of Dong, Palcoux and Plessis. This order separates the
extension to arbitrary complex fusion categories from the earlier
fifteen-ring result. The six non-self-dual rings are also recovered directly
using central induction and twist moments. The introduction records the
scope of Larson's pseudo-unitary result, the unitary result of
Edie-Michell--Izumi--Penneys, and the earlier sphericality and
low-rank classification results.

A putative nonspherical category must have self-dual simples, and its
Frobenius--Perron dimension $D$ must be quadratic or quartic and less than
$1000$. A shared convexity estimate replaces several separate quadratic
estimates. In the cubic case the argument stops once sphericality is proved.
A local calculation with the socle of the based order gives the inert-prime
obstruction. In the quartic case one square inequality treats the
ordinary and exceptional field cases. The text distinguishes the field of
a dimension character from the field of its codegree and explains where
their equality is needed. It also specifies the factor of two relating
codegrees in the sphericalization to the codegrees of its two quotients.
The short quadratic character table now states a necessary list, without
asserting that each row occurs.

The ordinary and signed enumerations now solve the same three linear
associativity equations, including the cases with zero coefficients.
Of the 1642 self-dual ordinary rings in the dimension window, 41 pass
the preliminary conditions. Their 1798 signed tuples give 623 distinct
pairs of an ordinary ring and a signed Casimir polynomial. Central norm
divisibility leaves two pairs, both with the same ordinary ring.
Their inverse-square codegree sums are $1$ and $23/24$, respectively.
Both violate the upper bound $7/8$. This proves sphericality without
using the pivotal classification. Invoking its fifteen-ring theorem at
this point removes the need to reproduce its pivotal case analysis.
The earlier global bank of 92795 signed orbits and its exclusion ledger
are no longer dependencies of the paper or its supplement.

\textbf{Checks.}
The complete finite replay regenerated all 1692 ordinary rings in the
window, including the 50 with a dual pair. Removing the dimension cut
gave 16065 rings, and the multiplicity-two brute-force check gave 27.
A separate direct enumeration produced the same 1798 signed tuples.
The arithmetic checks recovered all 14 short quadratic rows, verified
372 endpoint comparisons, and checked the two quartic-field estimates.
All fifteen model tensors and their positive dimension characters were
reconstructed and compared with the final table. The revised PDFs were
built from the sources included in the three submission archives.
Complete human verification of the mathematical arguments remains pending.

\medskip
Sébastien Palcoux\\
Beijing Institute of Mathematical Sciences and Applications
\end{document}
